\subsection{Network assembly and rhythm emergence}

Having established the circuit architecture, we now demonstrate how progressive integration of synaptic connections generates the complete swim CPG.
The four-cell network was constructed by sequentially activating synapses, revealing which interactions are necessary and sufficient for rhythm generation (Fig.~\ref{fig:assembly}).

In the initial configuration (Fig.~\ref{fig:assembly}A), the four interneurons are uncoupled.
Si3 neurons fire tonically at approximately 50~Hz with elevated excitability modeling excitatory drive from Si1 command neurons, while Si2 neurons remain quiescent.
When fast excitatory synapses from Si3 to contralateral Si2 are activated (Fig.~\ref{fig:assembly}B), Si2 neurons transition to tonic firing.
Addition of slow inhibitory feedback from Si2 to Si3 (Fig.~\ref{fig:assembly}C) produces spike-frequency modulation within the contralateral E/I loops (Si3R$\leftrightarrow$Si2L and Si3L$\leftrightarrow$Si2R), but consistent with experimental observations \cite{SakuraiKatz2016}, these loops alone do not generate autonomous bursting.

Reciprocal inhibition between contralateral Si3 neurons (Fig.~\ref{fig:assembly}D) immediately produces burst alternation, organizing Si3 activity into 2--3 second bursts separated by hyperpolarized interburst intervals.
Si2 neurons exhibit corresponding burst episodes driven by their Si3 partners.
This demonstrates that the Si3 half-center oscillator, when combined with the E/I modules, provides the phase-switching mechanism that converts continuous E/I dynamics into discrete alternating bursts.

The complete network (Fig.~\ref{fig:assembly}E), incorporating reciprocal Si2 inhibition and weak electrical coupling between E/I pairs, generates stable antiphase alternation with 3--4 second bursts that match experimental recordings.
Importantly, the stable rhythm is robust to the order in which synaptic connections are activated, confirming that network topology rather than assembly history determines the final dynamics.
This distributed architecture---where no single component is solely responsible for rhythm generation---predicts that the network should tolerate transient disruptions.
We test this prediction against experimental recordings.
